\section{Related Work}
\label{sec:related}

The motivation of CR-GNN is not only architectural but also domain-specific: biological graph learning often requires preserving local structural cues together with broader graph context. We therefore review the literature from six connected perspectives, starting with biomolecular and plant applications before moving to graph classification, message passing operators, deep-training challenges, hybrid architectures, and cross-layer aggregation.

\subsection{GNNs in Biomolecular and Plant Graph Learning}

Biological graph learning has moved well beyond generic graph benchmarks. In biomolecular settings, GNNs have already been used for structure-based protein function prediction \cite{gligorijevic2021structure}, large-scale protein function prediction across species \cite{you2021deepgraphgo}, and protein--protein interface pattern discovery \cite{reau2022deeprankgnn}. These studies show that graph models are practically relevant when the target depends on coordinated structural signals rather than on isolated features. The same pattern is now emerging in plant computational biology. Recent work applies graph-based modeling to plant vacuole protein identification \cite{suiIdentificationPlantVacuole2023}, plant biological-network analysis \cite{charulekhaUncoveringComplicatedPlant}, and plant functional annotation from knowledge graphs \cite{baongoImprovingPlantFunctional2025}. This literature does not validate our exact cross-residual design, but it makes the application setting concrete: in both biomolecular and plant tasks, the model often needs to preserve multi-scale biological information rather than only deepen a single message-passing chain.

Benchmark choice matters in this context. TU datasets such as PROTEINS, DD, and ENZYMES remain widely used because they are easy to reproduce and still expose meaningful differences in graph scale, sparsity, and classification difficulty \cite{morris2020tudataset}. For the present work, these datasets provide a stable protein-oriented evaluation core while keeping the narrative aligned with real biological graph-learning use cases.

\subsection{Graph Neural Networks for Graph Classification}

Graph-level prediction requires aggregating node-level features into fixed-size graph representations. Early graph neural networks \cite{gori2005new} laid the foundation for learning on graph-structured data through recursive message passing. The Message Passing Neural Networks (MPNN) framework \cite{gilmer2017neural} unified various GNN architectures under a common formalism, identifying message and vertex update functions as core components.

For graph classification, readout functions play a crucial role in converting node embeddings to graph-level representations. Simple approaches apply global pooling operations such as mean or sum pooling \cite{duvenaud2015convolutional}. However, these methods may not adequately capture hierarchical structural information. To address this, differentiable pooling methods have been proposed: DiffPool \cite{ying2018hierarchical} learns node cluster assignments to coarsen graphs hierarchically, SAGPool \cite{lee2019self} employs self-attention to identify salient nodes, and LocalPool \cite{cangea2018towards} enables sparse hierarchical classification. Recent surveys provide comprehensive overviews of graph pooling techniques \cite{liu2022graph,li2024graph}. While these approaches achieve strong performance, they primarily focus on pooling design rather than on how residual information should be reused across node-level and graph-level computation stages.

\subsection{Message Passing Architectures and Operators}

Various message passing operators have been developed, each imposing different inductive biases on neighborhood aggregation. Graph Convolutional Networks (GCN) \cite{kipf2016semi} approximate spectral convolutions through normalized adjacency matrices, performing low-pass filtering that smooths node features. GraphSAGE \cite{hamilton2017inductive} samples and aggregates features from local neighborhoods using mean, LSTM, or pooling aggregators, enabling inductive learning on unseen graphs. Graph Isomorphism Networks (GIN) \cite{xu2018powerful} use injective aggregation functions to achieve maximum discriminative power, matching the expressive power of the Weisfeiler-Lehman test.

Attention mechanisms have been widely incorporated to weight neighbor contributions adaptively. Graph Attention Networks (GAT) \cite{velivckovic2017graph} compute attention coefficients for each neighbor, allowing nodes to focus on relevant information. Personalized Graph Neural Networks (PGNN) \cite{zhang2019personalized} further extend attention mechanisms for personalized aggregation. Recent work generalizes Transformer architectures to graphs \cite{dwivedi2020generalization} and incorporates edge directionality \cite{beaini2020directional, rossi2023edge} for improved modeling of directed and heterophilic graphs.

Other operator variants include ARMA filters \cite{bianchi2020graph} for flexible frequency response, MixHop \cite{abu2019mixhop} for mixing multi-hop neighborhoods, and anti-symmetric convolutions \cite{gravina2022anti} for stable deep architectures. While these operators are typically studied and deployed individually within homogeneous architectures, a critical gap remains: the literature provides much less comparison of how different residual information-flow mechanisms behave across operator families under a unified training protocol.

\subsection{Deep GNN Training and Representation Degradation}

Training deep GNNs faces fundamental challenges due to representation degradation. The \textit{over-smoothing} phenomenon causes node representations to converge toward indistinguishable values as depth increases \cite{li2018deeper, chen2023revisiting, li2023non}. Theoretical analysis \cite{li2023non} provides non-asymptotic bounds on over-smoothing, revealing fundamental limitations in expressive power for deep architectures. Additionally, \textit{over-squashing} \cite{topping2021understanding, alon2020bottleneck} occurs when long-range dependencies cannot be effectively propagated due to graph bottlenecks, further constraining model performance.

To enable deep training, various techniques have been proposed. Residual connections \cite{he2016deep}, originally developed for computer vision, have been adapted to GNNs through residual gated graph convolutions \cite{bresson2017residual}. Normalization methods such as PairNorm \cite{zhao2020pair} and GraphNorm \cite{cai2021graphnorm} alleviate over-smoothing by normalizing node features, while learnable normalization \cite{li2022learning} adapts to different graph distributions. Regularization techniques like DropEdge \cite{rong2019dropedge} randomly drop edges during training to prevent overfitting and improve generalization.

Advanced architectures attempt to train extremely deep networks: techniques for training 1000-layer GNNs \cite{li2021training} and simplified deep GCN approaches \cite{oono2020simple} demonstrate that depth can improve performance with proper design. However, these methods primarily focus on single-path deepening and do not explicitly compare node-level residual reuse with graph-level residual propagation as two distinct ways of preserving information at depth.

\subsection{Multi-Operator and Hybrid Architectures}

Combining multiple aggregation schemes or architectural components has shown promise in enhancing representational capacity. Multi-hop attention \cite{wang2021multi} enables information flow across longer distances, while ARMA filters \cite{bianchi2020graph} provide flexible frequency responses beyond polynomial filters. Hybrid architectures incorporate edge directionality \cite{beaini2020directional, rossi2023edge} to capture directed relationships, and anti-symmetric convolutions \cite{gravina2022anti} improve stability in deep networks.

Surveys on attention-based GNNs \cite{sun2023attention} and graph embeddings \cite{cai2020understanding} provide comprehensive taxonomies of architectural variants. DenseGNN \cite{du2024densegnn} connects all layers densely, similar to DenseNet in computer vision, to maintain information flow. Hierarchical pooling methods \cite{pham2021hierarchical, ying2018hierarchical} create multi-scale representations through graph coarsening. However, most existing hybrid approaches emphasize component combination rather than a systematic comparison of where residual information should be injected and reused. Our work addresses this gap by contrasting node-level residual reuse, graph-level residual propagation, and cross-branch information exchange within a unified architecture family.

\subsection{Cross-Layer and Multi-Scale Aggregation}

Cross-layer connections aim to capture multi-scale information by aggregating features from different depths. Jumping Knowledge (JK) networks \cite{xu2018representation} adaptively select which layer's representation to use for each node, enabling flexible neighborhood range selection. Graph U-Net \cite{gao2019graph} adopts encoder-decoder architecture with skip connections between corresponding layers.

Propagation-based methods improve information flow across multiple hops: APPNP \cite{klicpera2018combining} combines personalized PageRank with neural predictions, while PairNorm \cite{zhao2020pair} and DropEdge \cite{rong2019dropedge} mitigate over-smoothing through normalization and regularization. Recent work on over-smoothing \cite{chen2023revisiting, zhao2020tackling} provides comprehensive analysis and solutions.

Despite these advances, most cross-layer methods operate within a single homogeneous branch using one type of operator. A fundamental limitation is that they rarely distinguish between residual reuse at the node-representation level and residual reuse at the graph-summary level. Our framework focuses on this distinction by introducing node-level and graph-level cross-residual mechanisms within a common graph classification design space. This perspective enables a more direct comparison of how different information-flow patterns affect robustness, depth tolerance, and graph-level discrimination in biological graph classification. With that context in place, we next formalize the shared graph-classification setting that makes the five-model comparison precise.
